\subsection{Properties of the Integers}

\begin{exercise}[0.2.1]
    For each of the following pairs of integers $a$ and $b$, determine their greatest common divisor, their least common multiple, and write their greatest common divisor in the form $ax + by$ for some integers $x$ and $y$.
    \begin{subexercises}
        \item $a = 20, b = 13$
        \item $a= 69, b = 372$
        \item $a = 792, b = 275$
        \item $a = 11391, b = 5673$
        \item $a= 1761, b = 1567$
        \item $a = 507885, b = 60808$
    \end{subexercises}
\end{exercise}

\begin{sol}
    \[
        \begin{array}{c|c|c|c}
            & (a, b) & \lcm(a, b) & ax + by \\
            \hline
            (a) & 1 & 260 & 2(20) - 3(13) \\
            (b) & 3 & 8556 & 27(69) - 5(372) \\
            (c) & 11 & 19800 & 8(792) - 23(275) \\
            (d) & 3 & 21540381 & 253(5673) - 126(11391) \\
            (e) & 1 & 2759487 & -105(1761) + 118(1567) \\
            (f) & 691 & 44693880 & 142(60808) - 17(507885)
        \end{array}
    \]
    Here, the $\lcm$ was computed using $ab = (a, b)\lcm(a, b)$.
\end{sol}

\begin{exercise}[0.2.2]
    Prove that if the integer $k$ divides the integers $a$ and $b$, then $k$ divides $as + bt$ for every pair of integers $s$ and $t$. 
\end{exercise}

\begin{sol}
    Since $k \mid a$ and $k \mid b$, there exist $x, y \in \zbb$ such that $a = kx$ and $b = ky$.
	Then for any $s, t \in \zbb$, we have
    \[
    	as + bt = (kx)s + (ky)t = k(xs + yt),
    \]
    which shows that $k \mid (as + bt)$.
\end{sol}

\begin{exercise}[0.2.3]
    Prove that if $n$ is composite, then there are integers $a$ and $b$ such that $n$ divides $ab$, but $n$ does not divide either $a$ or $b$.
\end{exercise}

\begin{sol}
    Since $n$ is composite, put $n = ab$ with $1 < a, b < n$.
	Then $n \mid ab$, but $n \nmid a$ and $n \nmid b$.
\end{sol}

\begin{exercise}[0.2.4]
    Let $a, b$ and $N$ be fixed integers with $a$ and $b$ nonzero, and let $d = (a, b)$ be the greatest common divisor of $a$ and $b$.
	Suppose $x_0$ and $y_0$ are particular solutions to $ax + by = N$.
	Prove, for any integer $t$, that the integers
    \[
    	x = x_0 + \frac bdt \quad \text{and} \quad y = y_0 - \frac adt
    \]
    are also solutions to $ax + by = N$.

    \note{This is, in fact, the general solution.}

    \mnote{Such equations are called \emph{Diophantine equations} and are an object of study in elementary number theory.}
\end{exercise}

\begin{sol}
    \begin{align*}
        ax + by & = a\left(x_0 + \frac bdt\right) + b\left(y_0 - \frac adt\right) \\
        & = ax_0 + \frac{ab}{d}t + by_0 - \frac{ab}{d}t \\
        & = ax_0 + by_0 = N \qh
    \end{align*}
\end{sol}

\begin{exercise}[0.2.5]
    Determine the value $\phi(n)$ for each integer $n \leq 30$, where $\phi$ denotes the Euler $\phi$-function.
\end{exercise}

\begin{sol}
    \[
        \begin{array}{c|c|c|c|c|c|c|c|c|c|c|c|c|c|c|c}
            n & 1 & 2 & 3 & 4 & 5 & 6 & 7 & 8 & 9 & 10 & 11 & 12 & 13 & 14 & 15 \\
            \hline
            \phi(n) & 1 & 1 & 2 & 2 & 4 & 2 & 6 & 4 & 6 & 4 & 10 & 4 & 12 & 6 & 8 \\
        \end{array}
    \]
    \[
        \begin{array}{c|c|c|c|c|c|c|c|c|c|c|c|c|c|c|c}
            n & 16 & 17 & 18 & 19 & 20 & 21 & 22 & 23 & 24 & 25 & 26 & 27 & 28 & 29 & 30 \\
            \hline
            \phi(n) & 8 & 16 & 6 & 18 & 8 & 12 & 10 & 22 & 8 & 20 & 12 & 18 & 12 & 28 & 8 \\
        \end{array} \qh
    \]
\end{sol}

\begin{exercise}[0.2.6]
    Prove the Well Ordering Property of $\zbb$ by induction, and prove the minimal element is unique.
\end{exercise}

\begin{sol}
    Suppose $S$ has no minimal element, and let $T$ be $\zp \setminus S$, the set of nonnegative integers not in $S$.
	Since $0$ is minimal in $\zp$, then $0 \notin S$, hence $0 \in T$.
	Fix $k \in \zp$, and suppose every $j$ such that $0 \leq j \leq k$ is in $T$.
	Observe $k + 1 \notin S$, for otherwise it would be minimal in $S$, since every integer less than or equal to $k$ is in $T$.
	Then $k + 1 \in T$.
	By induction, every nonnegative integer is in $T$, contradicting that $S$ is nonempty.
	Therefore, $S$ must have a minimal element.

    If $s_1, s_2 \in S$ are both minimal, then $s_1 \leq s_2$ and $s_2 \leq s_1$ imply $s_1 = s_2$.
\end{sol}

\begin{exercise}[0.2.7]
    If $p$ is a prime, prove that there do not exist nonzero integers $a$ and $b$ such that $a^2 = pb^2$.
\end{exercise}

\begin{sol}
    Suppose there existed $a, b$ such that $a^2 = pb^2$.
	If $(a, b) \neq 1$, divide both $a$ and $b$ by $(a, b)$ to get relatively prime integers.
	Then $p \mid a^2$, which implies $p \mid a$.
	Put $a = pk$ for $k \in \zbb$.
	Then
    \[
    	a^2 = (pk)^2 = p^2k^2 = pb^2,
    \]
    hence $b^2 = pk^2$.
	But then $p \mid b$, contradicting that $(a, b) = 1$.
\end{sol}

\begin{spex}[0.2.8]
    Let $p$ be a prime, $n \in \zp$.
	Find a formula for the largest power of $p$ which divides $n! = n(n - 1)(n - 2) \ldots 2 \cdot 1$.

    \hint{It involves the greatest integer function.}
\end{spex}

\begin{sol}
    In the product $n!$, a factor of $p$ appears in each of the integers $p, 2p, \ldots, \floor{n/p}p$.
	Hence, there are $\floor{n/p}$ factors of $p$ in $n!$.
	Similarly, a factor of $p^2$ appears in each of the integers $p^2, 2p^2, \ldots, \floor{n/p^2}p^2$, so there are $\floor{n/p^2}$ factors of $p^2$ in $n!$.
	This contributes an additional factor of $p$ that was counted only once in the previous step.
	Continuing this process, we find that the largest power of $p$ which divides $n!$ is
    \[
    	\floor*{\frac np} + \floor*{\frac{n}{p^2}} + \cdots + \floor*{\frac{n}{p^m}} = \sum_{i=1}^{m} \floor*{\frac{n}{p^i}},
    \]
    where $m = \floor{\log_p n}$ is the largest integer such that $p^m \leq n$.
\end{sol}

\begin{exercise}[0.2.9]
    Write a computer program to determine the greatest common divisor $(a, b)$ of two integers $a$ and $b$ and to express $(a, b)$ in the form $ax + by$ for some integers $x$ and $y$.
\end{exercise}

\begin{sol}
    Please see \verb|euclidean_algorithm.py| for the solution.
\end{sol}

\begin{exercise}[0.2.10]
    Prove for any given positive integer $N$ there exist only finitely many integers $n$ with $\phi(n) = N$ where $\phi$ denotes Euler's $\phi$-function.
	Conclude in particular that $\phi$ tends to infinity as $n$ tends to infinity.
\end{exercise}

\begin{sol}
    Fix $N \in \zp$, and consider $n \in \zp$ such that $\phi(n) = N$.
	Let $n = p_1^{\alpha_1} p_2^{\alpha_2} \cdots p_k^{\alpha_k}$ be the prime factorization of $n$.
	Then
    \[
    	\phi(n) = \prod_{i=1}^k p_i^{\alpha_i - 1}(p_i - 1) = N.
    \]
    Since each factor $p_i^{\alpha_i - 1}(p_i - 1)$ divides $N$, we have $p_i - 1 \leq N$, so $p_i \leq N + 1$, bounding the number of possible primes.
	For the exponents, $p_i^{\alpha_i - 1} \leq N$ gives $\alpha_i \leq 1 + \log_2 N$ since $p_i \geq 2$.
	With finitely many choices for each $p_i$ and $\alpha_i$, there are only finitely many $n$ with $\phi(n) = N$.

    If $\phi$ did not tend to infinity, there would exist some $M \in \zp$ such that $\phi(n) \leq M$ for infinitely many $n$.
	But then some value $N \leq M$ would satisfy $\phi(n) = N$ for infinitely many $n$, contradicting the above.
	Hence $\phi(n) \to \infty$ as $n \to \infty$.
\end{sol}

\begin{exercise}[0.2.11]
    Prove that if $d$ divides $n$ then $\phi(d)$ divides $\phi(n)$ where $\phi$ denotes Euler's $\phi$-function.
\end{exercise}

\begin{sol}
    Let $n \in \zp$ with $d \mid n$.
	Put their prime factorizations as
    \[
    	n = p_1^{\alpha_1} p_2^{\alpha_2} \cdots p_k^{\alpha_k}, \quad d = p_1^{\beta_1} p_2^{\beta_2} \cdots p_j^{\beta_j},
    \]
    where $j \leq k$ and $1 \leq \beta_i \leq \alpha_i$ for all $1 \leq i \leq j$.
	Then $\phi(d) \mid \phi(n)$, since
    \[
    	\frac{\phi(n)}{\phi(d)} = \prod_{i = 1}^j \frac{p_i^{\alpha_i - 1}(p_i - 1)}{p_i^{\beta_i - 1}(p_i - 1)} \cdot \prod_{i = j+1}^k p_i^{\alpha_i - 1}(p_i - 1) = \prod_{i = 1}^j p_i^{\alpha_i - \beta_i} \cdot \prod_{i = j+1}^k p_i^{\alpha_i - 1}(p_i - 1) \in \zbb. \qh
    \]
\end{sol}